\label{chap:83-02}

\section{Realizaciones geométricas del TRIT: regímenes elíptico, parabólico e hiperbólico y doble orientación temporal}

El TRIT conserva en un mismo estado el residuo publicado, el cociclo de
acarreo, la orientación y la memoria de composición. Sus \(3\) regímenes
algebraicos no son capas editoriales: son realizaciones correlativas del
operador cuadrático \(J_{\tau}^{2}=-\tau I\), enlazadas con las \(2\)
orientaciones del tiempo y con el estado de umbral.

\begin{HMTScientificOwnerSubordinated}
\def\HMTTRITFormal{1}
\def\HMTTRITDevelopments{1}
\input{sections/hmt/02_trit_geometrias}
\let\HMTTRITFormal\undefined
\let\HMTTRITDevelopments\undefined
\end{HMTScientificOwnerSubordinated}

% RESULTADO_RECUPERADO. Owner integro: sections/hmt/09b_algebras_cuadraticas.tex
% Destino causal: capitulo 2; realizaciones cuadraticas y generadores TRIT.
\begin{HMTScientificOwnerSubordinated}
\input{sections/hmt/09b_algebras_cuadraticas}
\end{HMTScientificOwnerSubordinated}

% RESULTADO_RECUPERADO. Owner integro: sections/hmt/03k_genealogia_operaciones_rev2.tex
% Destino causal: capitulo 2; operaciones levantadas y memoria del TPK.
\begin{HMTScientificOwnerSubordinated}
\input{sections/hmt/03k_genealogia_operaciones_rev2}
\end{HMTScientificOwnerSubordinated}

\section[APP, TRIT y TPK: soporte, orientación y transporte]
{APP, TRIT y TPK: soporte, orientación y transporte}

\subsection{APP como aritmética de fibras}

La APP no es una tabla ilustrativa: es el soporte aritmético donde se hace
visible la diferencia entre agregar y componer. En
\(
X_9=(\Znine)^2
\)
la suma acumula desplazamientos e historias alternativas; el producto compone
etapas y, cuando un factor no es invertible, pliega entradas distintas sobre
un mismo residuo. El levantamiento residuo--cociente despliega esas fibras y
conserva la genealogía contraída por la publicación modular.

Sea \(q(a,b)=a+b\pmod 9\). El dato publicado \(q(a,b)\) no determina el par
\((a,b)\). La fibra
\[
 q^{-1}(r)=\{(a,b):a+b\equiv r\pmod9\}
\]
es memoria de las alternativas. Para la multiplicación la contracción es aún
más marcada: las clases no unitarias \((0,3,6)\) producen fibras de tamaños
distintos. La APP publica, por tanto, una geometría de operaciones antes de
cualquier interpretación física.

El bloque central activo es
\[
B_c=\begin{pmatrix}7&2\\2&7\end{pmatrix}
   =7I+2S=9P_++5P_-,
\qquad
P_\pm=\frac{I\pm S}{2}.
\]
Su salto espectral local es \(9-5=4\); el coeficiente de acoplamiento es \(2\).
El valor \(54\) pertenece a un despliegue posterior y, en la fibra electrónica,
a \(D_1(\Gamma_e)\). Igualarlos por aparecer en una cadena numérica borraría
sus tipos.

\subsection{Geometría local de Lorentz y plano de Minkowski \texorpdfstring{\(1+1\)}{1+1}}

La normalización
\[
M_c=\frac{B_c}{\sqrt{45}}
=\exp\!\left(\frac12\log\frac95\,S\right)
\]
realiza un elemento de \(SO^+(1,1)\). APP contiene así, antes de la física
continua, una descomposición en modos \(P_+,P_-\), una brecha, un acoplamiento
y una rapidez. El resultado no es todavía el espacio-tiempo de Minkowski: es
la estructura algebraica que hará posible una representación lorentziana.

\subsection{TRIT como orientación y memoria local}

TRIT introduce el signo que la publicación modular omite. En una suma triádica
\[
 a+b=r+3c,
\]
el residuo \(r\) publica el resultado y el acarreo \(c\) conserva cómo se
obtuvo. Los estados \((+1,0,-1)\) no deben interpretarse sólo como \(3\)
colores. En la parametrización temporal actúan, respectivamente, como
retorno con memoria del pasado, umbral presente y apertura bidireccional al futuro;
en la realización algebraica distinguen contraflujo, fijación y flujo; en la
realización geométrica etiquetan las ramas hiperbólica, parabólica y elíptica
de una familia cuadrática. La compatibilidad entre estas lecturas exige mapas,
no una identificación verbal.

\begin{definition}[Estado TRIT enriquecido]
Un estado TRIT es un par \((r,c)\) acompañado de orientación. El primer
componente es publicable; el segundo pertenece a la fibra de memoria. El
cambio de lector puede olvidar \(c\), pero el transporte TPK no puede hacerlo.
\end{definition}

Esta separación permite formalizar presente, pasado y futuro sin convertirlos
en tres instantes absolutos. El presente es el lugar de acoplamiento de dos
orientaciones; pasado y futuro son sentidos de transporte. La inversión de
ruta intercambia ambos sentidos y deja el punto de acoplamiento como sección
de frontera.

\subsection{TPK como sistema de transporte}

TPK reúne soporte, orientación y memoria. Sus objetos son secciones preparadas
finitas o coinductivas; sus morfismos son rutas registradas compatibles con los
lectores. Cada ruta contiene
\[
\Gamma=(\text{origen},\text{destino},\text{hoja},\text{orientación},
\text{acarreo},\text{frontera},\text{registro}).
\]
La composición concatena caminos y compone memorias. La inversión revierte el
orden y la orientación. Los lectores, el calendario de \(108\) fases, la frontera y el
registro son operadores internos; no observadores externos añadidos.

Una operación publicada \(f:X\to Y\) admite un levantamiento
\[
 \widetilde f:X\longrightarrow Y\times M_f,
 \qquad \pi_Y\circ\widetilde f=f.
\]
Dados levantamientos fijados de \(f\) y \(g\), la composición tiene un
levantamiento canónico asociado
\[
 \widetilde{g\circ f}:X\longrightarrow Z\times M_f\times M_g.
\]
La memoria no es un comentario de la operación: es parte de su codominio
enriquecido.

\begin{proposition}[Genealogía de las operaciones]
En TPK, suma, producto, potencia, raíz, logaritmo, derivación e integración
admiten una lectura común:
\begin{center}
\begin{tabular}{@{}p{.18\textwidth}p{.31\textwidth}p{.38\textwidth}@{}}
\toprule
Operación & Acción & Memoria preservada \\
\midrule
Suma & agregación de alternativas & procedencia y acarreo \\
Producto & composición de etapas & orden y fibra plegada \\
Potencia & iteración en monoide & profundidad y órbita \\
Raíz & inversión multivaluada & rama y orientación \\
Logaritmo & coordenada de iteración & generador y calibre \\
Derivación & diferencia por escala & borde y orientación \\
Integración & composición de incrementos & camino y condiciones de frontera \\
\bottomrule
\end{tabular}
\end{center}
\end{proposition}

La potencia no introduce una operación ontológicamente desconectada de suma y
producto: itera la composición. La raíz invierte esa iteración y debe conservar
la selección de rama. El logaritmo registra profundidad. Esta jerarquía explica
por qué las potencias nonádicas y sus vacancias son informativas para APP: la
operación elevada conserva capas de memoria que una raíz digital aislada
contrae.

\subsection{Periodicidad temporal de (108) fases y leyes de retorno a (27), (54) y (108) iteraciones}

El calendario observable posee tres umbrales distintos. Tras 27 iteraciones
retorna la posición y se invierte la orientación; tras 54 se restauran posición
y orientación; tras 108 se restaura la fase observable completa. La memoria
del semigrupo positivo, sin embargo, no se borra. De forma esquemática,
\[
F_{\rm obs}^{27}(x,o)=(x,-o),\qquad
F_{\rm obs}^{54}(x,o)=(x,o),\qquad
F_{\rm obs}^{108}(x,o,\phi)=(x,o,\phi).
\]
Este orden jerárquico será crucial para separar propagación de ida, retorno,
orientación y ciclo completo. Un periodo observable finito puede coexistir con
una memoria genealógica no periódica.

\begin{falsadorBox}[title={Controles de tipo}]
No es admisible: identificar TPK con una lista de 34 operadores; llamar
«categoría fibrada» a la estructura sin levantamientos cartesianos; situar el
acarreo simultáneamente como coborde celular y como 2-cociclo sin separar sus
complejos; o convertir el bloque \(SO^+(1,1)\) en espacio-tiempo físico antes
de construir la representación métrica.
\end{falsadorBox}


\section{Fibración de Hopf, círculos de Villarceau, doble recubrimiento y holonomía de Möbius}

\subsection{Estructura compleja y fibración de Hopf}

Sea \(S^3\subset\mathbb C^2\) y
\[
h(z_0,z_1)=
\bigl(2\Re(z_0\overline z_1),
2\Im(z_0\overline z_1),|z_0|^2-|z_1|^2\bigr)
\in S^2.
\]
En el toro de Hopf
\[
T_\eta=\{(z_0,z_1):|z_0|=\cos\eta,
|z_1|=\sin\eta\},
\]
las curvas
\[
\gamma_+(t)=(\cos\eta\,e^{it},\sin\eta\,e^{it}),
\]
\[
\gamma_-(t)=(\cos\eta\,e^{it},\sin\eta\,e^{-it})
\]
poseen comportamientos distintos: \(\gamma_+\) es una fibra de Hopf y
\(\gamma_-\) se proyecta con grado dos. Bajo proyección estereográfica, las
clases homológicas \((1,1)\) y \((1,-1)\) producen las dos familias diagonales
de Villarceau.

La conjugación compleja total
\[
K(z_0,z_1)=(\overline z_0,\overline z_1)
\]
preserva cada familia no orientada y revierte su orientación. No intercambia
las dos familias. La conjugación parcial puede intercambiarlas, pero no es una
aplicación compleja lineal ni el antiunitario estándar global. Esta corrección
es esencial para no fabricar una simetría partícula--antipartícula a partir del
dibujo toroidal.

\subsection{Doble recubrimiento y retorno spinorial}

La aplicación \(SU(2)\to SO(3)\) es un doble recubrimiento. Una vuelta de
\(2\pi\) en \(SO(3)\) puede levantar a \(-I\) en \(SU(2)\); la vuelta de
\(4\pi\) recupera \(I\). La estructura de signo no convierte un círculo en
banda de Möbius. Para obtener una banda se necesita un fibrado real de línea
cuyo carácter de holonomía sea \(-1\).

Sea \(L_\chi\to S^1\) el fibrado asociado a
\[
\chi:\pi_1(S^1)\longrightarrow\{\pm1\},
\qquad \chi(1)=-1.
\]
Su espacio total es una banda de Möbius. La ruta \(C_M=-\operatorname{rev}\)
aporta el candidato de signo; el círculo Villarceau aporta el lazo; el
entrelazador físico debe probar que el primero es la holonomía del segundo.

\subsection{Descomposición espectral del operador orientado de memoria y derivación de la razón (5{:}2)}

Con \(B_c=7I+2S\), defínase el operador orientado de memoria
\[
B_{\rm mem}^{\varepsilon}=5I+2\varepsilon S,
\qquad\varepsilon=\pm1.
\]
Para \(\varepsilon=+1\),
\[
\Spec(B_{\rm mem}^{+})=\{7,3\},
\qquad\det B_{\rm mem}^{+}=21.
\]
Interpretar los valores \((7,3)\) como bordes ecuatoriales exterior e interior
en una carta positiva común produce
\[
R+r=7\ell,\qquad R-r=3\ell,
\]
y por tanto
\[
R=5\ell,\qquad r=2\ell,\qquad
\sin\beta=\frac rR=\frac25.
\]
El ángulo Villarceau es
\[
\beta=\arcsin\frac25
=23.5781784782\ldots^\circ.
\]

El teorema algebraico \(\{7,3\}\leftrightarrow(5,2)\) es exacto. La premisa
geométrica que interpreta el espectro como los dos bordes de un toro es una
interfaz declarada. No debe ocultarse ni emplearse el ángulo como blanco.

La realización del bloque completo \(B_c\) conserva otra lectura:
\(R:r=7:2\) y \(\beta=\arcsin(2/7)\). No son valores rivales; pertenecen a
dos funtores: bloque total y modo de memoria. La prioridad física del segundo
debe decidirse por el entrelazador con la dinámica, no por cercanía decimal con
otra constante.

\subsection{Coordenadas angulares conjugadas y normalización de la unidad de fase en la fibración de Hopf}

El ángulo nativo de HMT es una fase \(u\in\mathbb R/\mathbb Z\). Los lectores
publican
\[
\theta_{\rm rad}=2\pi u,
\qquad
\theta_{\rm deg}=360u.
\]
Radianes y grados son cartas distintas del mismo elemento circular. El radián
no es más «ontológico» por ser adimensional; su naturalidad procede de la
longitud de arco. Los grados proporcionan otra coordenada absoluta de la
vuelta. Cualquier relación con \(\alpha^{-1}\), \(\varphi\) o CKM debe
formularse primero como cociente o holonomía invariante y sólo después
evaluarse en una carta angular.

La identidad exacta ya disponible
\[
\tanh^2(2\log\varphi)=\frac59
\]
conecta el canal Perron--Fibonacci con el cociente espectral de \(B_c\). Su
complemento es \(4/9\), el salto normalizado. Esta relación es estructural.
No implica que todo ángulo construido con \(\varphi\) sea un ángulo físico.

\subsection{Brouwer, Möbius y grado}

El teorema del punto fijo de Brouwer, las transformaciones fraccionarias de
Möbius y la banda de Möbius son tres objetos distintos. Se relacionan sólo
después de fijar dominio, acción y fibrado. Las órbitas elípticas de adelanto y
retardo pueden modelarse como dos secciones de un fibrado con holonomía de
signo; la transformación de Möbius describe la acción conformal en la carta;
el teorema de Brouwer garantiza un punto fijo para mapas continuos del disco.
Ninguna de estas frases sustituye el diagrama de realización.

\begin{fronteraBox}[title={Resultado y frontera}]
Se demuestran el grado doble de la curva Hopf, la orientación de la
conjugación total y la construcción algebraica \(7,3\mapsto5,2\). La lectura
toroidal (5{:}2), la holonomía Möbius y la interpretación
partícula--antipartícula forman una cadena de interfaces que debe compartir un
único entrelazador. El volumen no fuerza ese entrelazador mediante una
coincidencia angular.
\end{fronteraBox}


\section{Realización hilbertiana finita del sistema de caminos y sus proyectores condicionales}

\subsection{Realización matricial de cada profundidad}

Para \(d\ge1\), definimos
\[
  \mathcal C_9^{(d)}=(\mathbb Z/9\mathbb Z)^d,
  \qquad
  H_d=\ell^2(\mathcal C_9^{(d)}),
  \qquad
  A_d=B(H_d)\cong M_{9^d}(\mathbb C).
\]
La base canónica \(\{\delta_x:x\in\mathcal C_9^{(d)}\}\) conserva la
etiqueta discreta de cada estado. La medida uniforme induce la traza
normalizada
\[
  \tau_d(a)=\frac{1}{9^d}\Tr(a).
\]

La factorización
\[
  \mathcal C_9^{(d+1)}
  =
  \mathcal C_9^{(d)}\times\mathbb Z/9\mathbb Z
\]
produce una identificación unitaria
\[
  H_{d+1}\cong H_d\otimes\mathbb C^9.
\]
Ésta es la entrada exacta de la torre operatorial. No se identifica por ello
\((\mathbb Z/9)^3\) con \(\mathbb F_3^6\) como grupos: comparten \(729\)
elementos, pero sus operaciones y cociclos de acarreo son distintos.

\subsection{Representación de Weyl–Heisenberg de la fase y el desplazamiento sobre la fibra trítica}

Sea \(\omega=e^{2\pi i/9}\). Sobre
\(H_1=\ell^2(\mathbb Z/9\mathbb Z)\), definimos
\[
  (Xf)(r)=f(r-1),
  \qquad
  (Zf)(r)=\omega^r f(r).
\]
Entonces
\[
  X^9=Z^9=I,
  \qquad
  ZX=\omega XZ.
\]
El grupo generado por \(X,Z,\omega I\) es el grupo de Heisenberg finito de
orden \(9^3=729\).

\begin{theorem}[Stone--von Neumann finito]
Toda representación unitaria irreducible del grupo de Heisenberg finito en
la que el centro actúa mediante el carácter primitivo
\(\omega I\) es unitariamente equivalente a la representación anterior.
\end{theorem}

\begin{proof}
Sea \(v\) un autovector de \(Z\). Los vectores
\[
  v,Xv,\ldots,X^8v
\]
tienen autovalores distintos de \(Z\), pues
\(ZX^kv=\omega^kX^kZv\). Son ortogonales y generan un subespacio
invariante. Por irreducibilidad forman una base de toda la representación.
En esa base, \(X\) es el desplazamiento cíclico y \(Z\) la fase diagonal.
\end{proof}

\begin{controlbox}
La unicidad falla si no se fija el carácter central primitivo o si se
permiten representaciones reducibles. Este teorema finito no debe confundirse
con la versión continua para las relaciones canónicas de conmutación.
\end{controlbox}

\subsection{Proyectores condicionales e incompatibilidad pre-Heisenberg}

Las hojas aditiva y multiplicativa de \(\APP\) definen subespacios de
información y sus proyectores ortogonales
\(P_{\Sigma,d},P_{\Pi,d}\). Las identidades construidas satisfacen
\[
 \|[P_{\Sigma,2},P_{\Pi,2}]\|
 =\frac{\sqrt2}{3},
 \qquad
 \|[P_{\Sigma,3},P_{\Pi,3}]\|
 =\frac{\sqrt3}{4}.
\]
La no conmutatividad no se deduce de una analogía cuántica: es un teorema
matricial finito. Su lectura «pre-Heisenberg» es posterior a ese cálculo.

\begin{proposition}
Los proyectores anteriores no pertenecen conjuntamente a una misma
subálgebra booleana de proyecciones.
\end{proposition}

\begin{proof}
En una subálgebra abeliana todos los proyectores conmutan. Los conmutadores
mostrados son no nulos.
\end{proof}

\subsection{Estructura hilbertiana de la publicación matricial y memoria genealógica preservada}

El espacio de Hilbert proporciona superposición lineal, proyección ortogonal
y cálculo espectral. HMT aporta el índice genealógico de la base y distingue
qué operador representa una transición, una medida o una sombra. La
realización no sustituye el \(\TPK\): lo representa.

\begin{sourcebox}
Los proyectores APP y el par Weyl--Heisenberg forman parte de la construcción
precedente. La organización \(H_d,A_d\) y la prueba de unicidad se establecen
aquí mediante la aplicación del teorema clásico correspondiente a los objetos
HMT.
\end{sourcebox}


\section{Estructura simpléctica finita de APP y fibración nonádica de observables}
\label{p12-83-c2-s4:gmc:chap:simplectica}

La geometría simpléctica no forma parte de los teoremas publicados sobre
cuadrados mágicos de De las Cuevas--Netzer. El vínculo legítimo pasa por un
corredor vecino y estándar: operadores de Weyl, espacio de fases discreto y
bases mutuamente insesgadas \citep{gibbons2004phase,planat2007rings}. En HMT
ese corredor adquiere una elevación nonádica y una memoria de fibra.

\subsection{La forma alternante y la conmutación de Weyl}

Sobre $V_3=\mathbb F_3^2$ definimos
\[
 \langle(q,p),(q',p')\rangle_{\rm sp}=qp'-pq'.
 \tag{GMC-5.1}
\]
Si $X|k\rangle=|k+1\rangle$,
$Z|k\rangle=\omega^k|k\rangle$ y $\omega=e^{2\pi i/3}$, los operadores
$W_{q,p}=X^qZ^p$ satisfacen, con estas convenciones,
\[
 W_uW_v=\omega^{-\langle u,v\rangle_{\rm sp}}W_vW_u.
 \tag{GMC-5.2}
\]
Por tanto, dos operadores conmutan exactamente cuando sus etiquetas son
ortogonales para la forma alternante. En dimensión dos, las subrectas
isotrópicas maximales son los cuatro puntos de
\[
 \mathbb P^1(\mathbb F_3)=\{[1:0],[0:1],[1:1],[1:2]\}.
 \tag{GMC-5.3}
\]
Cada una genera un contexto abeliano maximal de $M_3(\mathbb C)$; sus
proyectores espectrales constituyen una \PVM. Contextos asociados a líneas
distintas producen las cuatro bases MUB de \textup{(GMC-2.13)}.

\begin{proposition}[Descomposición qutrit por direcciones]
\label{p12-83-c2-s4:gmc:prop:descomposicion-masa-qutrit}
Sean $\mathcal D_a$ las cuatro subálgebras diagonales determinadas por las
líneas de \textup{(GMC-5.3)} y $\mathcal D_a^0$ sus partes de traza cero. Entonces
\[
 M_3(\mathbb C)
 =\mathbb CI\oplus
   \bigoplus_{a\in\mathbb P^1(\mathbb F_3)}\mathcal D_a^0
 \tag{GMC-5.4}
\]
como suma ortogonal para el producto de Hilbert--Schmidt. La cuenta de
dimensiones es $9=1+4\cdot2$.
\end{proposition}

\begin{proof}
Cada $\mathcal D_a^0$ tiene dimensión dos. El insesgamiento implica
$\Tr(A^*B)=0$ para $A\in\mathcal D_a^0$ y
$B\in\mathcal D_b^0$, $a\ne b$. Todas son ortogonales a $I$. Las
dimensiones suman nueve, la dimensión de $M_3(\mathbb C)$.
\end{proof}

Esta identidad explica por qué las doce proyecciones no son doce piezas
independientes arbitrarias: son cuatro resoluciones de la unidad cuyas partes
centradas descomponen todo el espacio operatorial qutrit.

\subsection{La elevación \texorpdfstring{$12\to4\times3$}{12 a 4 por 3}}

Sobre el módulo nonádico $V_9=(\mathbb Z/9\mathbb Z)^2$ se conserva la forma
alternante
\[
 \langle(a,b),(c,d)\rangle_9=ad-bc\pmod9.
 \tag{GMC-5.5}
\]
La línea proyectiva $\mathbb P^1(\mathbb Z/9\mathbb Z)$ está formada por
los vectores primitivos, módulo multiplicación por una unidad. Puede
representarse como
\[
 \{[1:t]:t\in\mathbb Z/9\mathbb Z\}
 \ \sqcup\ 
 \{[3u:1]:u\in\mathbb Z/3\mathbb Z\},
 \tag{GMC-5.6}
\]
y tiene $9+3=12$ puntos. La reducción módulo tres da el diagrama
\[
 \mathbb P^1(\mathbb Z/9\mathbb Z)
 \xrightarrow{\ r\ }
 \mathbb P^1(\mathbb F_3),
 \qquad |r^{-1}(a)|=3.
 \tag{GMC-5.7}
\]

La base $a$ de un contexto qutrit y el resultado $k$ de su \PVM también
forman un conjunto $4\times3$. Una \emph{calibración} es una biyección
\[
 \Xi:\mathbb P^1(\mathbb Z/9\mathbb Z)
 \longrightarrow
 \{P_{a,k}:a\in\mathbb P^1(\mathbb F_3),\ k\in\mathbb F_3\}
 \tag{GMC-5.8}
\]
que envía cada fibra de $r$ a los tres resultados del contexto situado sobre
$a$.

\begin{controlbox}
La fibración es canónica; la asignación del orden interno de cada fibra a los
tres resultados $k$ es un calibre. Una biyección no es todavía una
equivarianza. Para promover $\Xi$ a entrelazador debe demostrarse que las
acciones lineales, proyectivas y antiunitarias pertinentes se transportan en
todas las fibras. El calibre preliminar falla una comprobación antiunitaria en una
fibra; ese fallo se conserva como falsador, no se oculta.
\end{controlbox}

\subsection{Profundidad, acarreo y memoria sobre el espacio de fases discreto}

La geometría de Weyl sobre $\mathbb F_3^2$ explica conmutación, contextos y
MUB. HMT añade tres niveles que no se recuperan de la forma simpléctica por
sí sola:

\begin{enumerate}
 \item la elevación de una dirección ternaria a tres puntos de
 $\mathbb P^1(\mathbb Z/9\mathbb Z)$;
 \item la hoja, el acarreo y la orientación que distinguen estados con la
 misma reducción;
 \item la ruta TPK que transporta esas coordenadas y puede regresar a la
 fase sin regresar al estado.
\end{enumerate}

La aportación potencial a la literatura qutrit no es «otra construcción de
cuatro MUB», ya conocida. Es una geometría de \emph{raíces de contexto}: cada
contexto ternario posee tres antecedentes nonádicos, y esos antecedentes
pueden llevar una dinámica y una monodromía. El cuadrado de orden doce del
\Cref{p12-83-c22-s3:gmc:thm:qms-12-3} es la primera sombra normalizada de ese atlas.

\subsection{Separación tipada entre el solapamiento qutrit
\texorpdfstring{$1/3$}{1/3} y el lector HMT
\texorpdfstring{$1/\pi$}{1/pi}}

En una pareja de bases MUB qutrit,
\[
 \Tr(P_{a,k}P_{b,\ell})=\frac13.
 \tag{GMC-5.9}
\]
En el lector HMT aquí definido, la relación entre una hoja real y una
lectura volumétrica puede adoptar el cociente
\[
 \frac{\pi}{\pi^2}=\frac1\pi.
 \tag{GMC-5.10}
\]
Las dos expresiones no son intercambiables: $1/3$ es un solapamiento entre
proyectores en dimensión tres; $1/\pi$ es un peso escalar de otro lector.
La diferencia no impide coordinarlas. Fijada una \PVM
$\{P_0,P_1,P_2\}$, definimos una medición de nueve efectos:
\[
 F_k^{(\pi)}=\frac3\pi P_k\quad(k=0,1,2),
 \qquad
 F_r^{(\pi)}=\frac{1-3/\pi}{6}I_3\quad(r=3,\ldots,8).
 \tag{GMC-5.11}
\]
Es una \POVM porque $3/\pi<1$ y
\[
 \sum_{r=0}^{8}F_r^{(\pi)}
 =\frac3\pi I_3+\left(1-\frac3\pi\right)I_3=I_3.
 \tag{GMC-5.12}
\]
Si $Q$ es un proyector de rango uno de un contexto MUB distinto,
\[
 \Tr(QF_k^{(\pi)})
 =\frac3\pi\Tr(QP_k)
 =\frac3\pi\cdot\frac13
 =\frac1\pi.
 \tag{GMC-5.13}
\]

\begin{bridgebox}
La fórmula \textup{(GMC-5.11)} transforma exactamente la probabilidad MUB $1/3$ en
el lector $1/\pi$ mediante el calibre $3/\pi$, sin identificar ambos números.
Los seis efectos residuales conservan la unidad y completan una \POVM de
nueve resultados, de la que se obtiene otro \QMS cíclico $(9,3)$. Esta
construcción no identifica todavía la procedencia geométrica de la masa
residual; fija el tipo operatorial exacto que deberá tener.
\end{bridgebox}

\subsection{Grupo simpléctico y ecuación de equivarianza}

En dimensión dos sobre $\mathbb F_3$, el grupo simpléctico coincide con
$SL(2,\mathbb F_3)$ y permuta las cuatro líneas. Su acción se levanta,
proyectivamente, mediante el grupo de Clifford sobre las \PVM qutrit. Sea
$G_{\rm HMT}$ un grupo o grupoide de transformaciones del atlas nonádico y
$U_g$ el operador unitario o antiunitario candidato en la carta qutrit. La
condición que convierte una calibración en un puente dinámico es
\[
 \Xi(g\cdot\ell)=U_g\,\Xi(\ell)\,U_g^*,
 \qquad
 g\in G_{\rm HMT},\quad
 \ell\in\mathbb P^1(\mathbb Z/9\mathbb Z).
\tag{GMC-5.14}
\]
Si $g$ es una flecha y no sólo una simetría global, la ecuación debe incluir
origen y destino y se formula como transformación natural. La verificación de
\textup{(GMC-5.14)} es el paso que separa un etiquetado afortunado de una
transducción estructural.

\begin{sourcebox}
Las identidades simplécticas y MUB se apoyan en la literatura de fase finita
\citep{gibbons2004phase,planat2007rings}. No se atribuyen a los artículos de
cuadrados mágicos de De las Cuevas--Netzer. La fibración nonádica, la
calibración y los lectores de ruta constituyen la contribución propia de la
construcción HMT aquí estudiada.
\end{sourcebox}
